\documentclass[12pt,a4paper]{article}

\usepackage[margin=1in]{geometry}
\usepackage{setspace}
\usepackage{enumitem}
\usepackage[hidelinks]{hyperref}

\setstretch{1.15}

\begin{document}

\begin{center}
    {\LARGE \textbf{ENS 491 -- Graduation Project (Design)}}\\[1em]
    {\large (Draft) Proposal}
\end{center}

\vspace{1.5em}

\noindent \textbf{Project Title:} \hrulefill

\vspace{1em}

\noindent \textbf{Group Members:} \hrulefill

\vspace{1em}

\noindent \textbf{Supervisor(s):} \hrulefill

\vspace{1em}

\noindent \textbf{Date:} \hrulefill

\vspace{2em}

\noindent ENS 491--492 Proposal Report must contain the following sections (leave section titles unchanged).

\begin{itemize}[leftmargin=1.5em]
    \item You are required to write clearly and concisely.
    \item Each part of the proposal could further be improved and updated in your progress reports as you make further progress on your project. Project advisor may have additional requirements for the report.
\end{itemize}

\section{ABSTRACT}
Provide a brief but comprehensive summary of your project in approximately 250 words.

\section{INTRODUCTION}
This section should provide the general information on the project topic with the following specific details:

\begin{itemize}[leftmargin=1.5em]
    \item Explain the problem by giving the background information.
    \item Provide a literature survey on the existing studies/solutions (state of the art, commercial products, etc.) and state/elaborate the shortcomings or disadvantages of them.
    \item Mention if there are any patents related with your project.
    \item Justify the need for this project and state clearly the motivation behind the problem (your project).
\end{itemize}

\section{PROPOSED SOLUTION AND METHODS}
This part should introduce the general approach for the solution of the problem. You should address the following specific points:

\begin{itemize}[leftmargin=1.5em]
    \item What do you suggest to solve exactly?
    \item What is the scope (what is covered and not covered)?
    \item Explain why this project solves a ``complex problem''. In general, a problem is considered as complex if it does not have a straight-forward solution. Other characteristics of complex problems can be listed as follows:
    \begin{itemize}[leftmargin=1.5em]
        \item The problem involves wide-ranging or conflicting technical/scientific/engineering issues (multiple realistic constraints).
        \item The problem has no obvious solution and requires abstract thinking, originality in analysis to formulate suitable models.
        \item The problem requires research-based knowledge which allows a fundamentals-based, first principles analytical approach.
        \item The problem involves infrequently encountered issues.
        \item The problem has significant consequences in a range of contexts such as economic, environmental, social, manufacturability, safety, sustainability etc.
        \item The problem is a high level one involving various components.
    \end{itemize}
\end{itemize}

\subsection{Objectives/Tasks}
List the objectives/tasks of your project and corresponding intended results for each objective/task.

\subsection{Realistic Constraints}
Describe here the realistic constraints (such as economic, environmental, health-safety, manufacturability, etc.) that affected this project and how you addressed them.

\begin{itemize}[leftmargin=1.5em]
    \item \textbf{Economic:} Available budget (if no budget is available existing resources is the constraint); maintenance cost; prices of current similar products on market etc.
    \item \textbf{Environmental:} Pollution (air, noise, water etc.); global warming; energy saving; waste etc.
    \item \textbf{Social:} Government codes to protect society; workers’ rights; discrimination, personal privacy, etc.
    \item \textbf{Health and Safety:} Safety of workers; consumers and public; noise; hazardous materials and environment for workers; toxic materials; systems for increased equipment safety etc.
    \item \textbf{Manufacturability:} Designs with an impossibly small manufacturing tolerance, surface quality or mass; availability of chosen material; production of very hard materials; extremely small or thin sections etc. If the project includes software development, this item is mostly related to operational environment of the final product (such as computer, server and peripheral configurations, communication bandwidth, operating systems needed, etc.).
    \item \textbf{Sustainability (social, economic, and environmental):} Can the solution in your proposal survive under social, economic, and environmental aspects considered? A well-defined life span under the assumed normal operation (environmental) conditions; consideration of actual environmental factors (extreme working temperature, corrosive fluid, abrasive air, severe radiation in space, etc.) in design etc.
\end{itemize}

\subsection{Engineering/Scientific Standards}
Describe if your project had to satisfy specific engineering/scientific standards or codes of the field of specialty.

\begin{itemize}[leftmargin=1.5em]
    \item A technical or engineering standard is an established norm or requirement that are used to ensure quality, reliability, and safety.
    \item A code is a set of rules and specifications or systematic procedures that needs to be followed in the design, fabrication, installation, or inspection processes. The main goal of these codes is to set a bare/acceptable minimum level of safety.
\end{itemize}

\section{RISK MANAGEMENT}
All projects involve some risks. In this part of the report you need to discuss potential risks you may encounter during the course of your project development. You should also explain how you plan to address these risks if you encounter them. Some examples are given below.

\begin{itemize}[leftmargin=1.5em]
    \item What happens if you cannot acquire or use a planned resource?
    \item What happens if you require expertise you do not have in your team?
    \item What happens if you find out that one or more components in the project require longer delivery/development or manufacturing time?
    \item What is an alternative for the solution you proposed for each task (Plan B) in case the original one cannot be realized (contingency plan)?
\end{itemize}

\section{PROJECT SCHEDULE}
Provide a Gantt chart for the project.

\begin{itemize}[leftmargin=1.5em]
    \item Indicate the starting time and the required time period for completion of each task/objective.
    \item Indicate team member(s) responsible and working in each task; specify only one person as the task leader for each task.
    \item Discuss how delays in various tasks may affect the project and propose plans to eliminate or reduce their impact on the project outcome.
\end{itemize}

\section{ETHICAL ISSUES}
In this section, you are required to identify (if any) ethical aspects of your project/solution. For instance, are there possible legal/ethical consequences of the solution? Some examples are given below:

\begin{itemize}[leftmargin=1.5em]
    \item products implicitly using patent protected designs/concepts;
    \item materials that have better appearance but are toxic;
    \item under design for profit;
    \item products for secret survey of personal private life;
    \item solutions that may have negative affect on health.
\end{itemize}

\noindent If there aren’t any, you need to still keep this section and state that there are no ethical issues.

\section{REFERENCES}
Include all the references you used in your work. You can use any referencing style, but you should be consistent in formatting your references.

\end{document}